\section{CONSIDERAÇÕES FINAIS}

\subsection{Síntese e atendimento aos objetivos}

Este trabalho apresenta o desenvolvimento da Agendinha, concebida para tornar o agendamento digital simples e economicamente acessível a prestadores de serviços de renda limitada. A proposta substitui a negociação repetitiva de horários por mensagens e a consulta manual de compromissos por um link público de disponibilidade, preservando a autonomia do cliente e evitando cadastro ou instalação obrigatória.

Na versão documentada, o prestador pode criar uma conta, configurar o perfil e a disponibilidade, compartilhar seu link, consultar a agenda, criar bloqueios ou horários adicionais e cancelar reservas. O cliente pode escolher um horário, registrar a reserva sem conta e cancelar por meio de token próprio. Dessa forma, os fluxos implementados atendem aos objetivos funcionais definidos para o produto mínimo.

As decisões técnicas reforçam esse escopo. A confirmação da reserva ocorre em transação protegida contra sobreposição; RLS e validações limitam o acesso aos dados; o cache é tratado como otimização descartável; e as comunicações são registradas em outboxes para execução posterior. Esses mecanismos sustentam a integridade e a separação de responsabilidades, mas não substituem a avaliação da experiência com pessoas.

A simulação econômica indica que a mensalidade hipotética de R\$~5,00 pode cobrir a infraestrutura considerada no cenário estabelecido. A conclusão é condicional: demanda, consumo, tributos, suporte e demais custos ainda precisam ser observados antes de qualquer afirmação sobre sustentabilidade comercial.

A avaliação exploratória, com profissionais autônomos e participantes da área de software, permanece pendente. O protocolo e os campos de registro foram definidos na Seção~\ref{sec:avaliacao}, mas nenhum resultado de usabilidade, aceitação ou cobertura exploratória é afirmado nesta versão.

\subsection{Limitações}

\begin{itemize}
    \item a avaliação exploratória ainda não foi executada, portanto permanecem pendentes dados de conclusão de tarefas, dificuldades, defeitos e adequação ao cotidiano dos prestadores;
    \item os custos foram calculados com preços publicados, câmbio e demanda hipotéticos, sem incluir todas as despesas operacionais;
    \item o funcionamento depende de serviços externos de hospedagem, autenticação, cache e comunicação, cuja disponibilidade e cobrança podem variar;
    \item a versão atual é orientada à agenda de um prestador, sem gestão conjunta de vários profissionais no mesmo estabelecimento;
    \item não há processamento de pagamentos diretamente na plataforma.
\end{itemize}

\subsection{Trabalhos futuros}

A próxima etapa é executar o protocolo exploratório, consolidar os resultados nos campos reservados e priorizar ajustes a partir de dificuldades e defeitos reproduzíveis. A avaliação deverá preservar a distinção entre a percepção do público-alvo e a cobertura técnica realizada pelos profissionais de software.

No desenvolvimento funcional, duas evoluções são prioritárias. A primeira é permitir que \textbf{mais de um profissional cadastre e administre seus próprios horários em um mesmo estabelecimento}, com agendas independentes, escolha do profissional no fluxo de reserva e prevenção de conflitos de disponibilidade. A segunda é \textbf{possibilitar o pagamento pelo próprio sistema}, mediante integração com um provedor de pagamentos, contemplando confirmação de transações, estados de pagamento, cancelamentos e tratamento seguro das informações financeiras.

Como desdobramentos complementares, podem ser considerados reagendamento, lembretes automáticos e relatórios essenciais. A priorização dessas funcionalidades deverá considerar os resultados da avaliação exploratória e os custos reais de manutenção, preservando a simplicidade que orienta o projeto.

Conclui-se que a versão implementada oferece uma base coerente com a proposta de ampliar o acesso ao agendamento digital: escopo funcional reduzido, reserva autônoma pelo cliente e atenção ao custo mensal de operação. A confirmação de sua adequação ao público dependerá da avaliação planejada e da incorporação dos achados ao produto.

\FloatBarrier
